\section*{Capítulo \ref{ch:g11:infrared} --- La espectroscopia infrarroja}

\begin{solution}{exo:g11:infrared:1}
$\lambda = 1/3300 = \qty{3.03e-4}{cm} = \qty{3.03}{\micro m}$;
$\lambda = 1/1050 = \qty{9.52e-4}{cm} = \qty{9.52}{\micro m}$. La banda de
\qty{3300}{cm^{-1}}, de mayor \omterm{def:g11:infrared:wavenumber}{número de onda}, corresponde a la vibración
más rápida.
\end{solution}

\begin{solution}{exo:g11:infrared:2}
$\qty{10}{\micro m} = \qty{1.0e-3}{cm}$, así que $\sigma = 1/\num{1.0e-3} =
\qty{1000}{cm^{-1}}$.
\end{solution}

\begin{solution}{exo:g11:infrared:3}
El \ce{C=O} de una \omterm{def:g11:functional-groups:carbonyl}{cetona}; un \ce{O-H} de \omterm{def:g11:functional-groups:alcohol}{alcohol} unido por \omterm{def:g11:polarity-and-cohesion:hydrogen-bond}{enlaces de
hidrógeno}; los enlaces \ce{C-H} de una cadena carbonada.
\end{solution}

\begin{solution}{exo:g11:infrared:4}
La banda de \qty{1715}{cm^{-1}}: el enlace \ce{C=O}.
\end{solution}

\begin{solution}{exo:g11:infrared:5}
El instrumento mide la luz que atraviesa la muestra; donde un enlace
absorbe, pasa menos luz y la curva baja. $T =
\qty{100}{\percent}$ significa que no se absorbe nada de la luz de ese
\omterm{def:g11:infrared:wavenumber}{número de onda}.
\end{solution}

\begin{solution}{exo:g11:infrared:6}
El primero tiene una banda \ce{C=O} a \qty{1715}{cm^{-1}} y ninguna
banda \ce{O-H}: es una \omterm{def:g11:functional-groups:carbonyl}{cetona}, la butanona (llamada también
butan-2-ona), \ce{CH3-CO-CH2-CH3}. El segundo tiene una banda \ce{O-H} y
ningún \ce{C=O}: un \omterm{def:g11:functional-groups:alcohol}{alcohol}, por ejemplo el butan-1-ol.
\end{solution}

\begin{solution}{exo:g11:infrared:7}
Los grupos \ce{O-H} del ácido están unidos por \omterm{def:g11:polarity-and-cohesion:hydrogen-bond}{enlaces de hidrógeno} muy
fuertes (las \omterm{def:g7:atoms-and-molecules:molecule}{moléculas} se emparejan, cada \ce{O-H} unido al \ce{C=O} de
la otra), y cada \omterm{def:g7:atoms-and-molecules:molecule}{molécula} está unida de manera un poco distinta: la
banda se desplaza todavía más y se extiende por un intervalo muy ancho.
\end{solution}

\begin{solution}{exo:g11:infrared:8}
En el líquido (A), la banda \ce{O-H} es ancha, hacia
\qty{3350}{cm^{-1}}; en el gas (B), donde las \omterm{def:g7:atoms-and-molecules:molecule}{moléculas} están muy
separadas y no forman \omterm{def:g11:polarity-and-cohesion:hydrogen-bond}{enlaces de hidrógeno}, es una banda estrecha hacia
\qty{3666}{cm^{-1}}.
\end{solution}

\begin{solution}{exo:g11:infrared:9}
Difícilmente: hacia \qty{1730}{cm^{-1}} para el \omterm{def:g11:functional-groups:carbonyl}{aldehído} y
\qty{1715}{cm^{-1}} para la \omterm{def:g11:functional-groups:carbonyl}{cetona}, valores próximos. Lo decidiría la
comparación de todo el espectro, incluida la \omterm{def:g11:infrared:band}{región de la huella
dactilar}, con espectros de referencia, u otra técnica.
\end{solution}

\begin{solution}{exo:g11:infrared:10}
Un \omterm{def:g11:functional-groups:carboxylic-acid}{éster} (\ce{C=O} hacia \qty{1735}{cm^{-1}}, \ce{C-O} intenso, ningún
\ce{O-H}), por ejemplo el etanoato de etilo. El ácido etanoico es
\ce{C2H4O2}, no \ce{C4H8O2}. El ácido butanoico tiene la fórmula correcta,
pero presentaría la banda \ce{O-H} muy ancha de los ácidos: queda
descartado.
\end{solution}

\begin{solution}{exo:g11:infrared:11}
El propan-1-ol es un \omterm{def:g11:functional-groups:alcohol}{alcohol}: como el A. El ácido propanoico es un \omterm{def:g11:functional-groups:carboxylic-acid}{ácido
carboxílico}: como el D.
\end{solution}

\begin{solution}{exo:g11:infrared:12}
La banda \ce{O-H} ancha hacia \qty{3350}{cm^{-1}} disminuye mientras
crece una banda \ce{C=O} intensa hacia \qty{1715}{cm^{-1}}. La reacción
ha terminado cuando la banda \ce{O-H} ha desaparecido y la banda
\ce{C=O} ya no crece.
\end{solution}

\begin{solution}{exo:g11:infrared:13}
El ácido butanoico presenta una banda \ce{O-H} muy ancha de 2500 a
\qty{3300}{cm^{-1}}, y el \omterm{def:g11:functional-groups:carboxylic-acid}{éster}, ninguna; el \ce{C=O} del ácido está
hacia \qty{1710}{cm^{-1}}, más ancho, y el del \omterm{def:g11:functional-groups:carboxylic-acid}{éster} hacia
\qty{1735}{cm^{-1}}, estrecho, con una banda \ce{C-O} intensa hacia
\qty{1240}{cm^{-1}}. El ácido tiene un grupo \ce{O-H} unido por \omterm{def:g11:polarity-and-cohesion:hydrogen-bond}{enlaces
de hidrógeno}, que además afloja un poco su \ce{C=O}.
\end{solution}

\begin{solution}{exo:g11:infrared:14}
Etanol que haya quedado (el \omterm{def:g11:functional-groups:alcohol}{alcohol} con el que se fabrica el \omterm{def:g11:functional-groups:carboxylic-acid}{éster}) o
agua (del \omterm{def:g4:air-a-mixture-of-gases:air}{aire} o del lavado): los dos tienen grupos \ce{O-H} unidos por
\omterm{def:g11:polarity-and-cohesion:hydrogen-bond}{enlaces de hidrógeno}. Se puede comparar la muestra con un espectro de
referencia, secarla, destilarla o medir su temperatura de ebullición.
\end{solution}

\begin{solution}{exo:g11:infrared:15}
Un \omterm{def:g10:lewis-and-shape:multiple-bond}{enlace doble} es un muelle más rígido que un enlace sencillo entre los
mismos \omterm{def:g7:atoms-and-molecules:atom}{átomos}: vibra más deprisa, a un \omterm{def:g11:infrared:wavenumber}{número de onda} mayor (1715 frente
a \qty{1050}{cm^{-1}}). Los enlaces \ce{O-H}, \ce{N-H} y \ce{C-H}
implican todos un \omterm{def:g7:atoms-and-molecules:atom}{átomo} de hidrógeno, el más ligero de todos: unas bolas
ligeras unidas por un muelle vibran deprisa, de ahí los \omterm{def:g11:infrared:wavenumber}{números de onda}
de 2900--\qty{3600}{cm^{-1}}.
\end{solution}

\begin{solution}{pb:g11:infrared:1}
\textbf{1.} \ce{CH3-CH2-OH}; \ce{CH3-CO-CH3}; \ce{CH3-COOH}.

\textbf{2.} \ce{C2H6O}; \ce{C3H6O}; \ce{C2H4O2}.

\textbf{3.} Hidroxilo; carbonilo (\omterm{def:g11:functional-groups:carbonyl}{cetona}); carboxilo.

\textbf{4.} El etanol: \ce{O-H} y \ce{C-H}. La propanona: \ce{C-H} y
\ce{C=O}. El ácido etanoico: \ce{O-H}, \ce{C-H} y \ce{C=O}.

\textbf{5.} Ninguna banda \ce{C=O} y una banda \ce{O-H} ancha: el
etanol.

\textbf{6.} Una banda \ce{O-H} muy ancha de 2500 a \qty{3300}{cm^{-1}}
y un \ce{C=O} intenso hacia \qty{1710}{cm^{-1}}: el ácido etanoico.

\textbf{7.} La propanona: la banda \ce{C=O} muy intensa a
\qty{1715}{cm^{-1}}, sin \ce{O-H}.

\textbf{8.} Probablemente (vinagre, \omterm{def:g6:solutions-and-solubility:solute}{disolvente}, \omterm{def:g11:functional-groups:alcohol}{alcohol}), pero oler
líquidos desconocidos es peligroso; el espectro es seguro y concluyente.

\textbf{9.} Los grupos \ce{O-H} del ácido están unidos por \omterm{def:g11:polarity-and-cohesion:hydrogen-bond}{enlaces de
hidrógeno} más fuertes (\omterm{def:g7:atoms-and-molecules:molecule}{moléculas} emparejadas), así que la banda está más
baja y es mucho más ancha.

\textbf{10.} \ce{CH3-O-CH3}, \ce{C2H6O}; no tiene \ce{O-H}, así que no
hay banda \ce{O-H}.

\textbf{11.} La banda \ce{C=O}: hacia \qty{1730}{cm^{-1}} en el
propanal, y hacia \qty{1715}{cm^{-1}} en la propanona.

\textbf{12.} La banda \ce{O-H}: un \omterm{def:g11:functional-groups:carboxylic-acid}{éster} no tiene \ce{O-H}. Su \ce{C=O}
estaría hacia \qty{1735}{cm^{-1}}.

\textbf{13.} Los \omterm{def:g11:organic-skeletons:isomer}{isómeros} comparten la fórmula; sus \omterm{def:g11:functional-groups:functional-group}{grupos funcionales}
son distintos, y el espectro muestra los grupos.

\textbf{14.} \qty{1715}{cm^{-1}}.

\textbf{15.} $1715 \times 100 = \qty{1.715e5}{m^{-1}}$.

\textbf{16.} $\lambda = 1/\num{1.715e5} = \qty{5.83e-6}{m} =
\qty{5.83}{\micro m}$.

\textbf{17.} No: la luz visible va de unos 0.4 a \qty{0.75}{\micro m}.
Es radiación infrarroja.

\textbf{18.} $\num{1.05e5}\ \si{m^{-1}}$, así que $\lambda = \qty{9.52e-6}{m}
= \qty{9.52}{\micro m}$.

\textbf{19.} El enlace \ce{C=O} (mayor \omterm{def:g11:infrared:wavenumber}{número de onda}, menor longitud de
onda).

\textbf{20.} Unos \qty{5.8}{\micro m}.
\end{solution}
